% Осень26!B9; нумерация внутри столбца без пустых строк.
% Источник: лекция 7, §11 (строка 365)
\definitionticket{07}{Формулировка теоремы о неявном отображении в терминах систем уравнений}
\label{definition:07}

\begin{theorem}
Пусть $O\subset\mathbb R^{m+n}$ открыто, $F=(F_1,\ldots,F_n):O\to\mathbb R^n$ принадлежит классу $C^r$, $r\in\mathbb N\cup\{\infty\}$, и $(a,b)\in O$.
Рассмотрим систему
\[
  F_i(x_1,\ldots,x_m,y_1,\ldots,y_n)=0,
  \qquad i=1,\ldots,n.
\]
Пусть $F(a,b)=0$ и
\[
  \det F'_y(a,b)\ne0,
  \qquad F'_y=\left(\frac{\partial F_i}{\partial y_j}\right)_{i,j=1}^{n}.
\]
Тогда существуют открытые окрестности $P\subset\mathbb R^m$ точки $a$ и $Q\subset\mathbb R^n$ точки $b$, $P\times Q\subset O$, и единственное отображение $\varphi:P\to Q$ класса $C^r$ такие, что $\varphi(a)=b$ и
\[
  F(x,y)=0\quad\Longleftrightarrow\quad y=\varphi(x)
  \qquad ((x,y)\in P\times Q).
\]
Его производная задаётся формулой
\[
  \varphi'(x)=
  -\bigl(F'_y(x,\varphi(x))\bigr)^{-1}F'_x(x,\varphi(x)),
\]
где $F'_x=(\partial F_i/\partial x_j)_{i=1,\ldots,n;\,j=1,\ldots,m}$.
\end{theorem}
